%% kapa: exam=1 round=36 year=2025 period=1 form=A subject=real_estate subject_ko=부동산학원론 seq=13 number=93
%% source: Q-Net 감정평가사 기출문제 (cst00309) artlSeq=5231525 file=raw/1008/5231525/2025년 제36회 감정평가사 1차시험 1교시 문제지 원본.pdf
%% conversion: kordoc 4.24.51 -> tools/gen1.py (auto); review=proofread-sonnet; flags=none
\begin{question}{93}
용도지역ㆍ지구ㆍ구역제에 관한 설명으로 옳지 않은 것은?
\begin{choices}
  \item 용도지역ㆍ지구ㆍ구역제는 특정 토지를 용도지역 등으로 지정한 후 해당 토지를 이용 목적에 맞게 적용하는 제도이다.
  \item 용적률ㆍ건폐율 등의 밀도규제와 특정행위의 허가ㆍ불허가 등의 행위규제로 구성되어 있다.
  \item 용도지역은 토지를 경제적ㆍ효율적으로 이용하고 공공복리의 증진을 도모하기 위하여 서로 중복되지 아니하게 도시ㆍ군관리계획으로 결정하는 지역을 말한다.
  \item 용도지구는 용도지역의 제한을 강화 또는 완화하여 적용하며, 경관ㆍ안전 등을 도모하기 위하여 서로 중복되지 아니하게 도시ㆍ군관리계획으로 결정하는 지역을 말한다.
  \item 용도구역은 용도지역 및 용도지구의 제한을 강화 또는 완화하여 적용하며, 시가지의 무질서한 확산방지, 계획적이고 단계적인 토지이용의 도모 등을 위하여 도시ㆍ군관리계획으로 결정하는 지역을 말한다.
\end{choices}
\end{question}
\answer{4}
